\ifthenelse{\equal{\LanguageOption}{italian}}{%
	\thispagestyle{plain}
	\pdfbookmark[1]{Sommario}{sommario}
	\chapter*{Sommario}
	
	\guideinfo{Nella sezione \textit{Sommario}, presenta un riassunto conciso del tuo progetto, evidenziando i punti principali. Inizia con una breve dichiarazione del problema o dell'obiettivo, seguita da una descrizione del tuo approccio o metodologia. Riassumi i risultati o le scoperte principali, sottolineandone l'importanza o le implicazioni. Concludi con una o due frasi sul contributo complessivo o sull'impatto del tuo lavoro. Il riassunto deve essere chiaro e conciso, idealmente di 150-250 parole, affinché i lettori comprendano rapidamente la tua ricerca e la sua importanza.}
	
	\vspace{2em}
	\noindent\textbf{Parole chiave:} Parola chiave A, Parola chiave B, Parola chiave C.
}{%
	\thispagestyle{plain}
	\pdfbookmark[1]{Abstract}{abstract}
	\chapter*{Abstract}
	
	\guideinfo{In the \textit{Abstract} section, provide a concise summary of your project, highlighting the key points. Begin with a brief statement of the problem or objective, followed by a description of your approach or methodology. Summarise the main results or findings, emphasising their significance or implications. Conclude with a sentence or two on the overall contribution or impact of your work. Keep the abstract clear and focused, ideally within 150-250 words, to give readers a quick understanding of your research and its importance.}
	
	\vspace{2em}
	\noindent\textbf{Keywords:} Keyword A, Keyword B, Keyword C.
}

\MediaOptionLogicBlank